This appendix records the independent computational audit for
Theorem~\ref{thm:s5-classification}.  The theorem itself is proved in the
main text by finite classification and fiber counts.  The current
repository snapshot stores the expected audit output in
\[
\texttt{../03-certificates/s5\_expected\_output.txt}.
\]
The current replay script
\[
\texttt{../scripts/replay\_s5\_extension\_tiler\_audit.py}
\]
reproduces the aggregate S5 counts and the former rank and packing audit
certificates.  These computations are retained as independent audit
support; the proof of Theorem~\ref{thm:s5-classification} remains the
finite table plus the elementary obstructions in the main text.

The six-element counterexample in Theorem~\ref{thm:s6-counterexample} has
a separate self-contained verifier:
\[
\texttt{../scripts/verify\_s6\_p038\_biset\_counterexample.py}.
\]
It reconstructs the poset, the subgroups $H,K$, the multiplier set
$A=HK\sqcup H\tau K$, and checks that the $48$ translates of the
$15$-element tile cover all $720$ elements of $S_6$ exactly once.

The broader six-element laboratory status is recorded in
\[
\texttt{../03-certificates/s6\_classification\_status.md}
\]
and in the companion JSON file with the same base name.  Historical
Colab/RoundingSat artifacts for the five hard residual-clique cases are
kept as an audit trail only; the note
\[
\texttt{../03-certificates/legacy\_colab\_artifacts.md}
\]
marks them as timeout/UNKNOWN exploration rather than proof certificates.

\subsection*{Aggregate counts}

The script verifies:
\[
\begin{array}{lr}
\text{labeled posets} & 4231,\\
\text{unlabeled posets} & 63,\\
\text{series-parallel labeled posets} & 2791,\\
\text{series-parallel unlabeled posets} & 48,\\
\text{divisible labeled posets} & 3151,\\
\text{divisible unlabeled posets} & 51,\\
\text{non-series-parallel divisible labeled posets} & 360,\\
\text{non-series-parallel divisible unlabeled posets} & 3.
\end{array}
\]

\subsection*{The former rank obstructions}

For the first $|L(P)|=5$ type, one representative has transitive
relations
\[
\begin{gathered}
(0,1),(0,2),(0,3),(0,4),\\
(1,2),(1,3),(4,2).
\end{gathered}
\]
Its linear extensions are
\[
\begin{gathered}
(0,1,3,4,2),\quad (0,1,4,2,3),\quad
(0,1,4,3,2),\\
(0,4,1,2,3),\quad (0,4,1,3,2).
\end{gathered}
\]
The certificate is
\[
        \operatorname{rank}_{\mathbb F_5}(M)=115,\qquad
        \operatorname{rank}_{\mathbb F_5}(M\mid \mathbf 1)=116.
\]

For the second $|L(P)|=5$ type, one representative has transitive
relations
\[
(0,1),(0,2),(0,3),(1,2),(3,2),(4,1),(4,2).
\]
Its linear extensions are
\[
\begin{gathered}
(0,3,4,1,2),\quad (0,4,1,3,2),\quad
(0,4,3,1,2),\\
(4,0,1,3,2),\quad (4,0,3,1,2).
\end{gathered}
\]
The same rank certificate holds:
\[
        \operatorname{rank}_{\mathbb F_5}(M)=115,\qquad
        \operatorname{rank}_{\mathbb F_5}(M\mid \mathbf 1)=116.
\]

\subsection*{The former packing obstruction}

For the $|L(P)|=8$ type, one representative has transitive relations
\[
(0,1),(0,2),(0,3),(1,2),(4,2),(4,3).
\]
Its linear extensions are
\[
\begin{gathered}
(0,1,4,2,3),\quad (0,1,4,3,2),\quad
(0,4,1,2,3),\quad (0,4,1,3,2),\\
(0,4,3,1,2),\quad (4,0,1,2,3),\quad
(4,0,1,3,2),\quad (4,0,3,1,2).
\end{gathered}
\]
The translate family has $120$ distinct left translates.  Among the
translates disjoint from the base translate there are $96$ candidates,
and the maximum packing containing the base translate has size $14$.
Since a tiling would require $15$ translates, this type is not an
extension tiler.
